\section{Sample Interview Transcript}
\label{app:transcript}
\textbf{Interviewer:} How would you characterize the current state of carbon markets for small-scale renewable energy projects in sub-Saharan Africa?

\textbf{EXP-01:} Investors are still more interested in bigger, grid-connected projects where excess power can be sold back to the grid. Rural mini grids don't attract the same kind of commercial attention. Impact funds are still commercially focused, even if their LPs are a bit more patient. There's this shift happening, DFIs like the British have scaled down grants and shifted to more commercial investments. It's becoming more of a trade-not-aid mentality.

\textbf{Interviewer:} What types of renewable energy projects are successfully accessing carbon finance in the region?

\textbf{EXP-01:} Mini-grids just aren't seen as a viable asset compared to things like cookstoves. Intermediaries and registries don't prioritize them because the returns are smaller. If someone could aggregate projects under one umbrella and make it big enough, then maybe intermediaries would pay attention, because they could earn higher fees. But right now, they're focused on easier asset classes.

\textbf{Interviewer:} What are the primary barriers preventing mini-grid developers from accessing carbon markets?

\textbf{EXP-01:} The cost of entry is high, and fragmentation is a big issue. These developers are small and scattered, and the methodologies weren't really designed with them in mind. Verra and Gold Standard dominate most of the market, like 80\%, so it's really their policies that matter. From a policy perspective, it's their rules you have to influence.

\textbf{Interviewer:} How do transaction costs for mini-grid carbon certification compare to project revenues?

\textbf{EXP-01:} Covering the MRV costs and registration links straight back to investor returns. If the transaction costs are too high compared to what you can earn in credits, then the economics just don't work. That's why people keep trying to aggregate across countries, but it's tough to pull off.

\textbf{Interviewer:} At what project scale do carbon economics become viable for mini-grids?

\textbf{EXP-01:} I don't have a number, but megawatt-scale projects obviously do better than small rural systems. Maybe mini-grids tied to hospitals or schools that can connect to the main grid would make more sense. Mini-grid financing has better additionality prospects, you can argue that without carbon revenues, these projects wouldn't happen.

\textbf{Interviewer:} Have you worked with or observed mini-grid carbon projects in Africa? What were the key success factors or failure points?

\textbf{EXP-01:} I've seen people trying to aggregate projects across different countries. That's the holy grail. But it's hard because of transaction costs, small size, and policy fragmentation. Permanence and additionality still matter the same way whether it's a large grid or a mini-grid. The concepts don't change, but how you account for them does.

\textbf{Interviewer:} What role do intermediaries or aggregators play and why aren't there more serving the mini-grid sector?

\textbf{EXP-01:} They're not interested commercially. Mini-grids aren't attractive enough compared to other assets like cookstoves. If someone could bring a big aggregated bundle of projects, then maybe. But intermediaries would also need to invest in defining new methodologies, and right now that's not worth it to them.

\textbf{Interviewer:} Are current carbon methodologies well-suited to mini-grid contexts, or do they need adaptation?

\textbf{EXP-01:} Methodologies are evolving because of integrity concerns. There's been so much noise about greenwashing that buyers and registries are tightening up. Additionality is actually easier to prove with mini-grids, because without carbon revenue they often wouldn't happen. But the methodologies weren't designed with them in mind, so they'd need tweaking.

\textbf{Interviewer:} What are the biggest technical challenges for mini-grids in meeting MRV requirements?

\textbf{EXP-01:} MRV is expensive, and the administrative burden is high. For small projects it just eats too much of the value. People think methodologies will get better over time, but who influences the registries? It's not governments. It's mostly the private sector on the demand side.

\textbf{Interviewer:} How do voluntary carbon buyers view mini-grid credits compared to other project types? What's the pricing differential?

\textbf{EXP-01:} Buyers don't see mini-grids as a strong asset class, especially compared to cookstoves or REDD+. They just don't think the returns justify the effort. SSA will be a net seller, and the global north will be the buyers, that means the demand side decides what's acceptable. There's also more reputational risk now about integrity and greenwashing, so buyers are cautious.

\textbf{Additional insights documented from EXP-01:}
\begin{itemize}
	\item There's interest in using instruments like green bonds, insurance products, or SPACs to derisk and bundle assets.
	\item Article 6 rules and integrity debates are reshaping how additionality and permanence are handled.
	\item Some markets like Zimbabwe or South Africa could have more viable mini-grid carbon projects because they're so coal-reliant.
	\item Rural mini-grid scalability is still viewed skeptically compared to other clean energy assets.
	\item Registries and demand-side actors, not governments, are likely to influence how methodologies evolve.
\end{itemize}